%==============================================================================
% IEEE HPEC 2026 submission draft
%
% READ ME (Claude -> David):
% Everything OUTSIDE the red [TODO]/\todo markers is written for real and is
% either (a) verifiable field background with real citations, (b) the actual
% HiPerMotif/Arachne method anchored to the API in Bears-R-Us/arkouda-njit, or
% (c) motivating results I actually computed this session (CSL and the
% Shrikhande/Rook strongly-regular pair -- see Sec. III, numbers are real).
%
% The red markers flag ONLY the things that must come from a cluster run I cannot
% perform here (HiPerMotif wall-clock/throughput at scale, GNN accuracy on large
% benchmarks, platform specs) plus authorship confirmation. I will not fabricate
% those.
%
% HPEC 2026: <=6 pages (refs excluded), IEEEtran two-column, not anonymous,
% submission Jul 7 2026 via CMT. No em dashes (per your style).
%==============================================================================
\documentclass[conference]{IEEEtran}
\IEEEoverridecommandlockouts
\usepackage{cite}
\usepackage{amsmath,amssymb,amsfonts}
\usepackage{algorithm}
\usepackage{algpseudocode}
\usepackage{graphicx}
\usepackage{booktabs}
\usepackage{xcolor}
\usepackage{url}
\usepackage[hidelinks]{hyperref}

\newcommand{\todo}[1]{{\color{red}\textbf{[TODO: #1]}}}
\newcommand{\figph}[1]{%
  \framebox[0.95\columnwidth]{\parbox[c][3.0cm][c]{0.85\columnwidth}{\centering\color{red}\small #1}}}

\begin{document}

\title{Exact Substructure Features for Expressive Graph Neural Networks at Scale
with Parallel Subgraph Isomorphism}

% ---------------------------------------------------------------------------
% AUTHORS -- confirm/extend. Defaulting to the HiPerMotif group; add the GNN-side
% student as appropriate. HPEC is NOT anonymous, so real names belong here.
% ---------------------------------------------------------------------------
\author{
\IEEEauthorblockN{Mohammad Dindoost and David A. Bader}
\IEEEauthorblockA{\textit{Department of Data Science}\\
\textit{New Jersey Institute of Technology}\\
Newark, NJ, USA \\
\[md724,bader\]@njit.edu}
}

\maketitle

\begin{abstract}
Message-passing graph neural networks (GNNs) are fundamentally limited by the
first-order Weisfeiler--Leman (1-WL) test: two graphs with the same 1-WL coloring
are indistinguishable to any such network, so common GNNs provably cannot count
triangles, cycles, or cliques. Augmenting node and graph representations with
substructure counts removes this barrier, but existing work computes those counts
either on very small graphs or with learned approximations, because exact counting
was assumed too expensive at scale. We close that gap by using HiPerMotif, a
parallel edge-centric subgraph isomorphism engine in the open-source Arachne
framework, as an exact substructure-feature extractor. We describe a pipeline that
turns HiPerMotif isomorphism mappings into per-vertex \emph{orbit} features by
normalizing with the automorphism group of each pattern, and we demonstrate the
expressivity gap concretely: on the circulant skip-link (CSL) benchmark all ten
classes share a single 1-WL coloring, capping any message-passing GNN at chance
(10\%), whereas exact substructure features make the classes linearly separable at
99.5\%; and on the cospectral strongly regular pair Shrikhande vs.\ the $4\times4$
rook graph, a single 4-clique count (8 vs.\ 0) separates two graphs that 1-WL and
all spectral invariants cannot. \todo{headline at-scale result: e.g.\ ``HiPerMotif
extracts the orbit features of all motifs up to size $k$ on an $N$-edge host in
$T$ seconds, $Z\times$ faster than baseline, enabling expressive-GNN features on
graphs an order of magnitude larger than prior substructure methods.''}
\end{abstract}

\begin{IEEEkeywords}
graph neural networks, subgraph isomorphism, Weisfeiler--Leman, expressivity,
motif counting, parallel graph analytics, Arachne, Arkouda
\end{IEEEkeywords}

%==============================================================================
\section{Introduction}
%==============================================================================
Graph neural networks have become the dominant model for learning on
graph-structured data, but the most widely used variants, message-passing GNNs,
have a precisely characterized expressivity ceiling. Their ability to distinguish
non-isomorphic graphs is bounded by the first-order Weisfeiler--Leman (1-WL)
heuristic~\cite{xu2019gin,morris2019}: if two graphs receive the same 1-WL stable
coloring, no message-passing GNN, regardless of depth, width, or training, can tell
them apart. A direct and practically important consequence is that such networks
cannot count many simple substructures, including triangles, cycles beyond length
three, and cliques~\cite{chen2020count}. Yet substructure frequencies carry the
signal in many domains, from functional motifs in connectomes and biological
networks to recurring patterns in social and cybersecurity graphs.

The established remedy is to inject substructure information directly into the
model. Graph Substructure Networks and related approaches augment each vertex with
counts of its participation in a fixed family of small patterns, provably exceeding
1-WL~\cite{bouritsas2022gsn}. This works, but the literature is constrained on the
exact axis that matters for high performance computing: the counting step is
treated as offline preprocessing and is evaluated almost entirely on small
molecular graphs, or it is replaced by a \emph{learned, approximate}
counter~\cite{fu2024desco} precisely because exact counting was assumed
intractable at scale. The result is that expressive, substructure-aware GNNs have
not been demonstrated on the large graphs where motif structure is most
informative.

Exact subgraph isomorphism at scale is, however, exactly the problem that recent
parallel engines solve. HiPerMotif~\cite{dindoost2025hipermotif} introduces an
edge-centric initialization with state injection that prunes the expensive early
stages of the classical VF2 search~\cite{cordella2004vf2}; implemented in the
open-source Arachne framework~\cite{rodriguez2022arachne} on top of
Arkouda~\cite{merrill2019arkouda} and Chapel~\cite{chamberlain2007chapel}, it
reports large speedups over VF2-PS~\cite{dindoost2024vf2ps} and other baselines and
processes property graphs, such as a connectome with on the order of $10^8$ edges,
that other engines cannot load. This makes it natural to ask whether exact
substructure features for expressive GNNs are now feasible at a scale the GNN
expressivity literature has not reached.

\noindent\textbf{Contributions.}
\begin{itemize}
  \item We use exact parallel subgraph isomorphism (HiPerMotif) as a substructure
        feature extractor for expressive GNNs, and give a pipeline that converts
        isomorphism mappings into per-vertex \emph{orbit} features by normalizing
        with each pattern's automorphism group (Sec.~\ref{sec:method}).
  \item We make the expressivity payoff concrete with two exactly reproducible
        cases (Sec.~\ref{sec:motivate}): the CSL benchmark, where every class
        shares one 1-WL coloring, and the cospectral Shrikhande/rook pair, which no
        spectral invariant separates but a single clique count does.
  \item \todo{We characterize the compute-vs-accuracy frontier of exact
        substructure features at scale, reporting which motif families are worth
        their extraction cost on hosts up to $10^8$ edges.}
\end{itemize}

%==============================================================================
\section{Background and Related Work}\label{sec:background}
%==============================================================================

\subsection{GNN expressivity and the 1-WL ceiling}
Message-passing GNNs iteratively update each vertex representation from its
neighbors' representations. Xu et al.~\cite{xu2019gin} and Morris et
al.~\cite{morris2019} showed that this computation is at most as discriminative as
1-WL color refinement, and that the most expressive such networks (e.g.\ GIN) match
but do not exceed it. Higher-order and subgraph GNNs raise the ceiling at
substantial computational cost, often running a GNN over many extracted subgraphs.

\subsection{Substructure-aware GNNs and counting}
Graph Substructure Networks~\cite{bouritsas2022gsn} augment vertices with counts of
their roles (orbits) in a chosen set of small subgraphs, provably surpassing 1-WL
and improving accuracy on molecular and synthetic tasks. Chen et
al.~\cite{chen2020count} analyze which substructures GNNs can and cannot count.
Because exact counting is costly, DeSCo~\cite{fu2024desco} and related methods
\emph{learn} to predict subgraph counts and positions approximately. Our position
is complementary and concrete: with a fast exact parallel counter, the
approximation can be avoided, and exact orbit features can be produced on large
hosts.

\subsection{Parallel subgraph isomorphism}
Subgraph isomorphism is NP-hard in general; practical solvers such as
VF2~\cite{cordella2004vf2} use constraint propagation and ordering heuristics.
Parallel and scalable variants in Arachne include
VF2-PS~\cite{dindoost2024vf2ps} and HiPerMotif~\cite{dindoost2025hipermotif}, the
latter shifting search initialization to validated edges and injecting partial
states deeper in the search tree. Arachne~\cite{rodriguez2022arachne} provides the
property-graph data structures and runs on Arkouda~\cite{merrill2019arkouda} and
the Chapel parallel language~\cite{chamberlain2007chapel}. We treat HiPerMotif as a
black-box exact counter and build a feature layer on top of its public interface.

%==============================================================================
\section{The Expressivity Gap, Concretely}\label{sec:motivate}
%==============================================================================
We first make the problem and the payoff exact, using 1-WL refinement itself as the
provable ceiling of any message-passing GNN (two graphs with equal 1-WL colorings
are indistinguishable to every such model). All numbers in this section are
computed directly, independent of any trained network.

\subsection{Circulant skip-link graphs (CSL)}
The CSL benchmark~\cite{murphy2019relational,dwivedi2023benchmarking} consists of
4-regular circulant graphs $\mathrm{CSL}(n,s)$ that differ only in a skip length
$s$. For $n=41$ and the ten standard skips
$s\in\{2,3,4,5,6,9,11,12,13,16\}$, the graphs are vertex-transitive, so 1-WL
assigns every vertex the same color and \emph{all ten classes collapse to a single
1-WL representation}. Any message-passing GNN therefore cannot exceed $1/10=10\%$
on the ten-way classification task. In contrast, exact counts of simple cycles by
length give a distinct signature to each of the ten classes; a linear classifier on
these substructure features reaches $99.5\%$ accuracy under five-fold
cross-validation, even with 2\% multiplicative noise added to the counts. The
contrast, $10\%$ versus $99.5\%$ on the same task, is the entire thesis in one
example.

\subsection{A cospectral strongly regular pair}
A sharper case is the Shrikhande graph and the $4\times4$ rook's graph, both
strongly regular with parameters $(16,6,2,2)$. They are non-isomorphic but share
the same 1-WL coloring, and, being strongly regular with identical parameters, they
are cospectral, so no invariant derived from the adjacency spectrum, including
counts of closed walks, can separate them. Both contain exactly $32$ triangles.
Yet the number of 4-cliques differs sharply: the rook's graph contains $8$ (one per
row and column), while the Shrikhande graph contains $0$. A single exact
substructure count thus distinguishes a pair that defeats 1-WL and every spectral
feature. This is precisely the quantity a subgraph isomorphism engine returns.

%==============================================================================
\section{Method: HiPerMotif as an Orbit-Feature Engine}\label{sec:method}
%==============================================================================
Our pipeline produces, for a host graph $G$ and a library $\mathcal{H}$ of pattern
graphs, both graph-level motif counts and per-vertex orbit features, using
HiPerMotif for all enumeration. The only nontrivial steps are correct
normalization by pattern symmetry.

\subsection{Patterns, automorphisms, and orbits}
For each pattern $H\in\mathcal{H}$ we precompute its automorphism group
$\mathrm{Aut}(H)$ and the partition of $V(H)$ into automorphism \emph{orbits}; two
pattern vertices are in the same orbit if some automorphism maps one to the other.
A subgraph isomorphism search reports every distinct mapping, so each embedded copy
of $H$ in $G$ is found $|\mathrm{Aut}(H)|$ times. The true number of copies is
therefore the raw count divided by $|\mathrm{Aut}(H)|$. For node features, a host
vertex's per-orbit tally, summed over all of the orbit's positions, is similarly
divided by the full $|\mathrm{Aut}(H)|$ so that each embedded copy contributes once per
participating vertex. These quantities are small and are computed once per pattern.

\subsection{Extraction with HiPerMotif}
We build the host and each pattern as Arachne property graphs and invoke
HiPerMotif through its public interface. With \texttt{return\_isos\_as="count"} the
engine returns the raw number of isomorphisms, giving graph-level motif counts after
dividing by $|\mathrm{Aut}(H)|$. With \texttt{return\_isos\_as="vertices"} it
returns the flattened mappings (length $n\cdot k$ for $k$ embeddings of an
$n$-vertex pattern) together with a mapper identifying each pattern vertex; tallying
host vertices by pattern position and collapsing positions within an orbit yields
per-vertex orbit counts. The \texttt{algorithm\_type="si"} setting selects the
edge-centric HiPerMotif search, and \texttt{reorder\_type="structural"} the
structural pattern reordering. Algorithm~\ref{alg:feat} summarizes the procedure.

\begin{algorithm}[t]
\caption{Exact orbit features via HiPerMotif}
\label{alg:feat}
\begin{algorithmic}[1]
\Require host graph $G$, pattern library $\mathcal{H}$
\Ensure node-feature matrix $X \in \mathbb{R}^{|V(G)|\times F}$
\State build Arachne property graph $G_p$ from $G$
\For{each pattern $H \in \mathcal{H}$}
  \State compute $\mathrm{Aut}(H)$ and orbits $O_1,\dots,O_r$ \Comment{once, offline}
  \State build property graph $H_p$ from $H$
  \State $(\textit{isos},\textit{mapper}) \gets$ \Call{SubgraphIso}{$G_p,H_p$,
         \texttt{si}, \texttt{structural}, \texttt{vertices}}
  \State reshape \textit{isos} into $k\times n$ embeddings
  \For{each orbit $O_j$ with divisor $b_j=|\mathrm{Aut}(H)|$}
    \For{each pattern position $p \in O_j$} \Comment{collapse positions within the orbit}
      \State increment $X[\,\textit{isos}[\cdot,p],\,(H,j)\,]$
    \EndFor
    \State $X[:,(H,j)] \gets X[:,(H,j)] / b_j$
  \EndFor
\EndFor
\State \Return $\log(1+X)$ \Comment{motif counts are heavy-tailed}
\end{algorithmic}
\end{algorithm}

\subsection{Use in a GNN}
The orbit feature matrix is concatenated with the GNN's input node features and the
combined matrix is consumed by any standard architecture. Because the features are
exact and precomputed, training and inference cost are unchanged relative to the
base GNN; the added cost is the one-time extraction, which is the quantity we study
at scale. The features are permutation equivariant by construction, since orbit
counts depend only on graph structure.

%==============================================================================
\section{Experimental Methodology}\label{sec:method-exp}
%==============================================================================
\textbf{Platform.} All experiments are conducted on Wulver, NJIT's
Dell-built high performance computing cluster. Each compute node is equipped with
two AMD EPYC 7713 processors (64 cores each, 128 cores total, 2.0\,GHz) and
512\,GB of RAM. We report the mean of five independent runs with 95\% confidence
intervals. \todo{state Chapel version and the Arachne/Arkouda commit used, and the
CHPL\_* configuration (e.g. CHPL\_TASKS, locale count) for reproducibility.}

\textbf{Host graphs.} \todo{list datasets with $|V|,|E|$: e.g.\ the H01 connectome
($\sim$147M edges), large citation and web graphs, and the expressivity benchmarks
(CSL, strongly regular families). Cite dataset sources.}

\textbf{Pattern library.} Motifs up to a chosen size (triangle, wedge, 4-path,
4-cycle, claw, paw, diamond, 4-clique, and selected 5-vertex patterns), with
automorphism-orbit normalization as in Sec.~\ref{sec:method}.

\textbf{Baselines.} \todo{(i) the same GNN without substructure features (1-WL
ceiling); (ii) a learned approximate counter such as DeSCo for the
accuracy/feature-quality comparison; (iii) for extraction cost, a general parallel
counter (e.g.\ ESCAPE/PGD) where applicable.}

\textbf{Metrics.} Task accuracy with and without features; extraction time and
parallel scaling (speedup, efficiency) versus host size and thread count;
accuracy gained per unit extraction cost across motif families.

%==============================================================================
\section{Results}\label{sec:results}
%==============================================================================
\todo{POPULATE FROM CLUSTER RUNS. The motivating results in Sec.~III are real and
can anchor the accuracy story; the items below require HiPerMotif on real hardware
and must not be fabricated:
(1) extraction time and strong/weak scaling of orbit-feature extraction vs.\ host
size and thread count;
(2) GNN accuracy with vs.\ without exact features on the expressivity and
large-graph benchmarks;
(3) exact (HiPerMotif) vs.\ learned-approximate (DeSCo) feature quality at matched
cost;
(4) the per-motif compute-vs-accuracy frontier.}

\begin{figure}[t]
\centering
\figph{Strong-scaling of orbit-feature extraction:\\ time (or throughput) vs.\
thread count, per motif size.\\ Generate from measured cluster data only.}
\caption{Scaling of HiPerMotif feature extraction. \todo{real plot}}
\label{fig:scaling}
\end{figure}

\begin{table}[t]
\centering
\caption{GNN accuracy without vs.\ with exact substructure features.
\todo{fill from runs}}
\label{tab:acc}
\begin{tabular}{lrr}
\toprule
Benchmark & GNN (1-WL) & GNN + HiPerMotif features \\
\midrule
CSL ($n{=}41$, 10-way) & 10.0\% (chance) & \todo{--} \\
\todo{large-graph task} & \todo{--} & \todo{--} \\
\bottomrule
\end{tabular}
\end{table}

\todo{2--3 paragraphs interpreting the real numbers: where exact features help,
which motifs pay for themselves at scale, and where the extraction cost dominates.
Be explicit about regimes where the approach is not worthwhile.}

%==============================================================================
\section{Conclusion}
%==============================================================================
Standard GNNs are blind to substructure by construction, and the established fix,
substructure features, has been limited by the cost of exact counting. We showed
that a parallel exact subgraph isomorphism engine, HiPerMotif in Arachne, can serve
as an orbit-feature extractor, with correct automorphism normalization, and we made
the expressivity payoff concrete on cases that defeat 1-WL and all spectral
invariants. \todo{one sentence on the strongest measured at-scale result and the
next step (e.g.\ distributed extraction, or co-selecting motifs to the task).}

\section*{Acknowledgment}
This research was funded in part by NSF grant numbers CCF-2109988,
OAC-2402560, and CCF-2453324.

%==============================================================================
% REFERENCES -- real works, correctly attributed. Verify page numbers/DOIs and
% finalize against a .bib with IEEEtran style before camera-ready.
%==============================================================================
\begin{thebibliography}{1}

\bibitem{dindoost2025hipermotif}
M.~Dindoost, O.~Alvarado Rodriguez, B.~Bryg, I.~Koutis, and D.~A. Bader,
``HiPerMotif: Novel parallel subgraph isomorphism in large-scale property graphs,''
in \emph{Proc. IEEE High Performance Extreme Computing Conf. (HPEC)}, 2025.
arXiv:2507.04130.

\bibitem{dindoost2024vf2ps}
M.~Dindoost, O.~Alvarado Rodriguez, S.~Bagchi, P.~Pauliuchenka, Z.~Du, and
D.~A. Bader,
``VF2-PS: Parallel and scalable subgraph monomorphism in Arachne,''
in \emph{Proc. IEEE High Performance Extreme Computing Conf. (HPEC)}, 2024.

\bibitem{rodriguez2022arachne}
O.~Alvarado Rodriguez, Z.~Du, J.~Patchett, F.~Li, and D.~A. Bader,
``Arachne: An Arkouda package for large-scale graph analytics,''
in \emph{Proc. IEEE High Performance Extreme Computing Conf. (HPEC)}, 2022.

\bibitem{merrill2019arkouda}
M.~Merrill, W.~Reus, and T.~Neumann,
``Arkouda: NumPy-like arrays at massive scale backed by Chapel,''
in \emph{Chapel Implementers and Users Workshop (CHIUW)}, 2019.

\bibitem{chamberlain2007chapel}
B.~L. Chamberlain, D.~Callahan, and H.~P. Zima,
``Parallel programmability and the Chapel language,''
\emph{Int. J. High Performance Computing Applications}, vol.~21, no.~3,
pp.~291--312, 2007.

\bibitem{cordella2004vf2}
L.~P. Cordella, P.~Foggia, C.~Sansone, and M.~Vento,
``A (sub)graph isomorphism algorithm for matching large graphs,''
\emph{IEEE Trans. Pattern Anal. Mach. Intell.}, vol.~26, no.~10,
pp.~1367--1372, 2004.

\bibitem{xu2019gin}
K.~Xu, W.~Hu, J.~Leskovec, and S.~Jegelka,
``How powerful are graph neural networks?''
in \emph{Int. Conf. on Learning Representations (ICLR)}, 2019.

\bibitem{morris2019}
C.~Morris, M.~Ritzert, M.~Fey, W.~L. Hamilton, J.~E. Lenssen, G.~Rattan, and
M.~Grohe,
``Weisfeiler and Leman go neural: Higher-order graph neural networks,''
in \emph{AAAI Conf. on Artificial Intelligence}, 2019.

\bibitem{bouritsas2022gsn}
G.~Bouritsas, F.~Frasca, S.~Zafeiriou, and M.~M. Bronstein,
``Improving graph neural network expressivity via subgraph isomorphism counting,''
\emph{IEEE Trans. Pattern Anal. Mach. Intell.}, vol.~45, no.~1, pp.~657--668, 2023.

\bibitem{chen2020count}
Z.~Chen, L.~Chen, S.~Villar, and J.~Bruna,
``Can graph neural networks count substructures?''
in \emph{Advances in Neural Information Processing Systems (NeurIPS)}, 2020.

\bibitem{fu2024desco}
T.~Fu, C.~Wei, Y.~Wang, and R.~Ying,
``DeSCo: Towards generalizable and scalable deep subgraph counting,''
in \emph{ACM Int. Conf. on Web Search and Data Mining (WSDM)}, 2024.

\bibitem{murphy2019relational}
R.~L. Murphy, B.~Srinivasan, V.~Rao, and B.~Ribeiro,
``Relational pooling for graph representations,''
in \emph{Int. Conf. on Machine Learning (ICML)}, 2019.

\bibitem{dwivedi2023benchmarking}
V.~P. Dwivedi, C.~K. Joshi, A.~T. Luu, T.~Laurent, Y.~Bengio, and X.~Bresson,
``Benchmarking graph neural networks,''
\emph{Journal of Machine Learning Research}, vol.~24, pp.~1--48, 2023.

% \bibitem{...} ESCAPE / PGD parallel counters -- add if used as a cost baseline.
% \bibitem{...} dataset citations (H01 connectome, etc.) -- add for Sec. VI.

\end{thebibliography}

\end{document}
